\section{Fundamentação e trabalhos relacionados}
\label{sec:fundamentacao}

\subsection{Governança adaptativa em organizações pequenas}

A governança de TI pode ser operacionalizada por estruturas, processos e mecanismos relacionais que definem direitos decisórios, coordenam negócio e tecnologia e monitoram a geração de valor. Em organizações pequenas, a utilidade desses mecanismos depende de fatores contingenciais como porte, centralização, maturidade, setor e disponibilidade de competências. O estudo de caso de \citet{levstek2022adaptive} mostra que uma parcela relevante dos mecanismos analisados é situacional e adverte que modelos universais podem ser custosos e difíceis de aplicar em PMEs. Como se trata de um único caso, o resultado fundamenta o desenho adaptativo, mas não demonstra por si só qual configuração é superior em outros contextos.

O estudo exploratório de \citet{silva2018baseline}, baseado em entrevistas semiestruturadas, investiga mecanismos de linha de base para governança empresarial de TI em PMEs. O registro institucional consultado confirma essa abordagem; o número de participantes e a distribuição detalhada dos mecanismos não foram verificados no texto integral. No GEAR, comunicação entre negócio e TI, acompanhamento de orçamento e responsabilidades são escolhas de desenho explicitadas em rotinas, sem atribuir a esse estudo a validação dessas práticas locais.

A transformação digital acrescenta uma dimensão temporal ao problema. A revisão sistemática de \citet{sagala2024transformation} recomenda que PMEs considerem sua linha de base e suas limitações, com avanços incrementais e investimento contínuo em aprendizagem. Essa orientação sustenta dois requisitos do GEAR: a entrada por uma fase curta de estabilização e o aumento gradual de práticas conforme evidências de maturidade. A revisão também registra concentração geográfica de estudos, o que recomenda cautela ao transferir seus achados para micro e pequenas empresas brasileiras.

\subsection{Bases normativas e gerenciais}

O acervo histórico relaciona o ciclo local de avaliar, dirigir e monitorar à ISO/IEC~38500:2024 \citep{iso38500_2024}. O texto integral da norma não foi verificado nesta revisão; a associação histórica não demonstra conformidade. No GEAR, ADM-Lite nomeia uma adaptação local, distinta do Architecture Development Method do TOGAF. O COBIT~2019 amplia esse domínio ao oferecer objetivos, componentes e fatores de desenho para um sistema de governança ajustado ao contexto \citep{isaca2018cobit}. O GEAR não reproduz o catálogo do COBIT; usa sua orientação a objetivos e desenho como origem para um subconjunto rastreável de decisões, responsabilidades e métricas.

O ITIL~4 oferece a principal referência para execução e entrega de serviços. Seus princípios orientadores, que orientam adoção e adaptação às necessidades específicas, conforme a apresentação pública consultada \citep{peoplecert2026public}, são compatíveis com a proposta de adequação ao contexto. No GEAR, essa orientação é materializada por canal único de demanda, quadro Kanban com limite de trabalho em progresso, ciclos semanais e um PRD de uma a duas páginas para iniciativas que exijam desenvolvimento. Essa seleção não representa uma implementação integral do sistema de valor de serviço do ITIL.

Para segurança, o NIST CSF~2.0 organiza resultados de gestão de risco nas funções Governar, Identificar, Proteger, Detectar, Responder e Recuperar \citep{nist2024csf}. Os CIS Controls~v8 acrescentam salvaguardas priorizadas e grupos de implementação; o IG1 é apresentado como higiene cibernética essencial e reúne 56 salvaguardas \citep{cis2026ig1}. O GEAR usa essas referências para formar uma linha de base centrada em inventário de ativos críticos, identidade e privilégio mínimo, backup verificável e resposta a incidentes. A expressão ``NIST-Lite'' usada no artefato designa uma adaptação interna e não uma publicação ou certificação do NIST.

\begin{table}[htbp]
  \centering
  \caption{Posicionamento das referências e do GEAR. A comparação descreve escopo e mecanismo, não efetividade observada.}
  \label{tab:posicionamento}
  \small
  \begin{tabularx}{\textwidth}{@{}L{2.3cm}L{3.0cm}R R@{}}
    \toprule
    Referência & Escopo predominante & Força aproveitada & Tradução no GEAR \\
    \midrule
    ISO/IEC 38500 & Referência histórica de governança & Texto integral não verificado nesta revisão & Associação histórica ao ADM-Lite, sem alegação de conformidade \\
    COBIT 2019 & Governança e gestão de informação e tecnologia & Objetivos, componentes e fatores de desenho & Rastreabilidade, responsabilidades e portfólio mínimo \\
    ITIL 4 & Gestão de produtos e serviços digitais & Foco em valor, fluxo e melhoria iterativa & Canal único, Kanban, ciclos semanais e PRD curto \\
    NIST CSF 2.0 & Resultados para gestão de risco cibernético & Funções de governança e ciclo de risco & Organização do pilar de Segurança Crítica \\
    CIS Controls v8 & Salvaguardas cibernéticas priorizadas & IG1 e ações mensuráveis de higiene essencial & Checklist mínimo e evidências operacionais \\
    \bottomrule
  \end{tabularx}
\end{table}

\subsection{IA como assistência, não como fonte de verdade}

Há evidência de que sistemas de IA generativa podem elevar produtividade em tarefas delimitadas. No experimento controlado de \citet{peng2023copilot}, 95 desenvolvedores profissionais implementaram um servidor HTTP em JavaScript; o grupo com GitHub Copilot concluiu a tarefa 55,8\% mais rápido, com intervalo de confiança de 95\% entre 21\% e 89\%. Em outro contexto, \citet{brynjolfsson2025work} analisaram a introdução de assistência por IA entre 5.172 agentes de suporte e observaram aumento médio de 15\% nas questões resolvidas por hora, concentrado entre trabalhadores menos experientes. Esses resultados motivam a hipótese de aceleração documental, mas não podem ser transferidos diretamente para governança de TI nem usados como estimativa de ganho do GEAR.

A engenharia de prompts oferece meios para fornecer contexto, restrições e exemplos a modelos sem alterar seus parâmetros \citep{sahoo2025prompt}. Em pesquisa qualitativa, o mapeamento de \citet{barros2025qualitative} identifica potencial de automação, mas também dependência de prompts, imprecisões ocasionais e limitações contextuais. Por isso, o módulo de IA do GEAR exige dados de origem, indicação de lacunas, citações quando aplicáveis e aprovação humana antes que uma saída gere decisão ou ação.

Referências históricas sem fonte primária confirmada foram separadas dos fundamentos desta edição. O registro de revisão bibliográfica conserva as lacunas e as decisões de restrição.

\subsection{Lacuna e distinção do artefato}

Os trabalhos existentes sustentam a adaptação contingencial, os mecanismos mínimos e a implantação incremental, enquanto as normas oferecem cobertura conceitual robusta. Permanece, porém, uma lacuna entre essas orientações e uma implementação integrada que uma equipe muito pequena consiga iniciar, medir e revisar. O GEAR ocupa esse espaço ao combinar três domínios essenciais, uma porta de entrada temporal, artefatos compactos e uma camada opcional de automação com controles de validação.

A novidade reivindicada é de composição e operacionalização, não a criação isolada dos mecanismos de governança, agilidade ou segurança. O valor científico esperado está na justificativa explícita das escolhas, na rastreabilidade entre fonte e prática, e na avaliação do ajuste organizacional sob protocolo comparativo. Essa formulação torna verificável tanto a contribuição quanto seus limites.
